% Source: 19_Chapter_11_Emmanuel_Christ_In_the_Creeds.txt
\bookchapter[Emmanuel!]{CHAPTER 11}{EMMANUEL!}{Christ in the Creeds}

\leadin{I}{N THE FOOTHILL REGION} of Mount Hermon, Jesus once asked his disciples, “Who do men say
that I am?” They told him that most people associated him with the prophets of Israel. When he
pressed them further, “And how about you? Who do you say I am?” Peter responded, “You are the
Messiah, the Son of the living God.”

Men have given a thousand answers to Jesus’ question. Some have said, “He is an unusual Jewish rabbi
preaching a kingdom of love.” Others answer, “No, he is a social revolutionary whose primary purpose
is the overthrow of Rome’s tyrannical rule.” Still others claim, “He is a misguided dreamer who
looked for God to step into history and establish justice on the earth.”

Whatever the views of men, the church through the centuries has confessed with Peter, Jesus Christ
is the Messiah, the Son of the living God. He is more than a subject for Christian study: he is the
object of Christian devotion. Theologians call this mystery the incarnation, the enfleshment of God.
Hymn writers extol the merits of “Emmanuel,” meaning “God with us.”

During the Imperial Age of the church, when emperors pressed the pastors for creedal statements
accurately expressing the Christian faith, the church came to speak of the God-man. In 451 a general
council at Chalcedon, not far from Constantinople, affirmed that Jesus Christ was “complete in
Godhead and complete in manhood, truly God and truly man~.~.~. in two natures, without confusion,
without change, without division or without separation~.~.~. coming together to form one person.”

Most Roman Catholic, Eastern Orthodox, and Protestant Christians to this day consider that statement
orthodox Christian belief. Sadly, some Christians in Egypt and Syria and India do not. The reasons
for their disagreement lie in the fifth-century attempts to speak clearly about the Event. How did
the classic doctrine of the incarnation come to be formulated? What does it mean?

\section{THE WORD BECAME FLESH}

We call this area of theology Christology because it raises the question, “Who was Jesus Christ?”
What was the relation of the divine life and the human life in this unique person, the Christian
Savior?

The very existence of this question in the life of the church is profoundly significant. So far as I
know, Islam has no Mohammedology and Buddhism has no Buddhology. The debate in the history of
Christianity is a monument to the uniqueness of the One Christians call the Son of God.

The Imperial Age did not create the question of the incarnation; it simply debated it. The mystery
of the God-man was central to Christian worship long before it became central to Christian thinking.
“A deep instinct,” J. S. Whale once told the undergraduates at Cambridge University, “has always
told the Church that our safest eloquence concerning the mystery of Christ is in our praise. A
living Church is a worshiping, singing Church; not a school of people holding all the correct
doctrines.”

Whale meant that the most treasured hymns of the church have always treated Christ as an object of
worship. We find the beating heart of Christian experience not in the church’s creed but in its
music.

A major task remains for devout scholarship, however, because God in flesh is more than Christian
sentiment; he is Christian reality. The apostles spoke of “the Image of the Invisible God,” “the
Word made flesh,” and “the Lamb slain from the foundation of the world.” Second-and third-century
believers rejected Ebionism and Gnosticism because they were so clearly distortions of this truth.

The fourth- and fifth-century debates about the meaning of the incarnation were not aimed at an
explanation of Christ. These Christians knew that Jesus Christ defies explanation because he fits no
class. He is unique. The great merit of the creeds is that they left the mystery intact.

It is no surprise, however, that the human mind revolts against these creeds from Nicaea and
Chalcedon. Any professor who has tried to trace the Christological debates for first-year
seminarians has seen the reaction. But as Charles Williams once pointed out, it is a revolt of
“immature sensibility.” It is natural.

The church councils appear beyond feeling. As a result we have the immature and romantic devotions
to “the simple Jesus, the spiritual genius, the broad-minded international Jewish workingman, the
falling-sparrow and grass-of-the-field Jesus.” But, said Williams, these will not serve the faith
because the Christian idea from the beginning believed that in Jesus Christ earth and heaven meet,
both man and God.

Christians of an earlier day felt this reality deeply. They also knew that the obligation to be
intelligent is a moral obligation, so they searched for some statement that would correlate the
human life and the divine life in Jesus Christ.

In the early church two famous schools of theology offered contrasting emphases on the Word becoming
flesh (John 1:14). One of these was at Alexandria; the other was at Antioch. The Alexandrians
emphasized how the divine Word assumed human flesh (humanity in general). The Antiochians emphasized
how the divine Word was joined to the man named Jesus. The inherent danger or risk for the
Alexandrians was to so emphasize the union of the divine and human that the humanity was
overshadowed by the divinity of Jesus. The inherent danger for the Antiochians was to underemphasize
the union of the divine and human in Jesus to the point that his divinity was overshadowed by his
humanity.

The early, leading voice at Alexandria was Origen, who, in speaking of Jesus Christ, coined the term
God-man. Drawing upon ideas from the famous Greek philosopher Plato, Origen developed a deep,
fervent mysticism centering in the divine Word (Logos). His ideas focused on the thought that in
Christ the encounter between God and mankind had taken place perfectly and that Christians must
strive to imitate it.

Later Christians developed this mystical idea. One of the Cappadocian fathers, Gregory of Nyssa,
taught that in Christ, the Logos—the one divine person—had united in himself the divine and human
natures. Both natures existed by themselves and distinct from one another, yet they were not
separated but arranged in such a way that their attributes were mutually exchangeable.

The Antioch school of theologians normally interpreted the Scriptures in a more straightforward
historical manner. Major teachers of this position tended to stress the human figure of the Gospels.
They found saving virtue in Jesus’ example and achievement. In Christ the human will, which in other
men turns freely to sin, proved obedient and victorious.

Antiochene theologians consequently stressed the concrete manhood of Christ. As David F. Wright
explains it, the union of body and soul did not in any way affect the completeness and normality of
the human nature. After the Word became flesh, the two natures remained distinct. In Antiochene
teaching they could easily seem like two beings, God and man, Son of God and son of Mary, joined or
associated rather than personally united. As a vessel indwelt by the Word, Jesus was not unlike
prophets and apostles, except that he enjoyed perfect fullness of grace and power. According to one
theologian the Word lodged in the man Jesus as in a temple.

The debate over the meaning of the Event raged for generations in part because political influence
was at stake. After Christianity emerged as the official religion of the empire under Theodosius,
the structure of the church centered upon a few powerful figures. The bishops in the chief towns of
the imperial provinces came to be called archbishops. The term for the official center of a bishop’s
jurisdiction and authority is see. Those bishops in the premier cities of the empire—Rome,
Constantinople, Alexandria, and Antioch—were considered highest of all and were called patriarchs.
Throughout the fourth and fifth centuries these four powerful patriarchs were attempting to extend
the prestige and power of their spiritual offices.

Generally speaking, Alexandria and Rome tended to support each other, and Antioch to go along with
Constantinople. On Alexandria’s part this reflected the jealousy that proud and ancient city felt at
the rise of the upstart capital in the east, Constantinople. Rome, too, while for the time content
to extend its influence over the imperial dominions in the west, was not happy with the growing
arrogance of the “new Rome” in the east. On the other hand, Antioch and Alexandria had long been at
loggerheads as rivals in the east. If Antioch could not gain preeminence, it preferred to see it go
to the church in the new capital rather than to its old rival on the Nile.

\section{MAJOR HERESIES ABOUT CHRIST}

In this climate the Christological debate stretched over a century and was the primary passion in
the churches of the East. Between 350 and 450 heresies arose, each of them forcing the churches to
greater clarity in their answer to the question, “Who is Jesus Christ?”

The first position advanced and rejected was associated with a pastor of Laodicea named Apollinaris,
a younger friend of Athanasius. Reacting to the teaching from Antioch, Apollinaris struck upon the
idea of approaching the question from the view of what we would call psychology. He felt that human
nature embraced the body and the soul. At the incarnation, however, the divine Word (Logos),
according to Apollinaris, displaced the animating and rational soul in a human body, creating a
“unity of nature” between the Word and his body. Humanity, he felt, was the sphere, not the
instrument of salvation, so he could speak of “one enfleshed nature of the divine Word.” Here was
the Alexandrian stress on the deity of Christ, but only a body represented Christ’s human nature.

Objections to Apollinaris’s position quickly arose. Does not the gospel picture Jesus as a complete
and genuine human? And if the Word displaced the rational soul of human nature, with its powers of
choice and sin, how can man be fully redeemed? Without complete solidarity with humanity, how could
the Word secure the salvation of mankind? As Gregory of Nazianzus put it, “What has not been assumed
cannot be restored.”

In this atmosphere of criticism, the second general council of the church, meeting at Constantinople
in 381, effectively silenced the Apollinarian teaching. It simply was not an adequate description of
the incarnation.

The second heresy was associated with the name Nestorius, a famous preacher at Antioch before the
emperor, in 428, made him bishop of Constantinople. The imperial capital gave Nestorius a platform.
From it he tried to defend the position of his teacher in the faith, Theodore, bishop of Mopsuestia,
near Antioch. Like his instructor, Nestorius rejected a popular designation of Mary as the
“God-bearer, Mother of God.”

In rejecting the phrase, Nestorius made it appear that he held that Christ joined two persons. He
did not deny the deity of Christ; but in emphasizing the reality and integrity of the Savior’s
humanity, Nestorius pictured the relation between the two natures in terms of a moral conjunction or
a merging of wills rather than that of an essential union. Although he never divided Christ into two
sons, Son of God and son of Mary, he refused to attribute to the divine nature the human acts and
sufferings of the man Jesus. Once he said, “I hold the natures apart, but unite the worship.” He
insisted that calling Mary “Mother of God” was tantamount to declaring that the divine nature could
be born of a woman, or that God could be three days old.

As it turned out, Nestorius was condemned more for political than for doctrinal reasons. In the hot
seat of Constantinople, he was soon widely hated for his assaults on Jews and heretics. He speedily
incurred the hostility of Cyril, the patriarch of Alexandria (412–44). Cyril was a distinguished
preacher and theologian, but ruthless in debates. He was especially alarmed by Nestorius’s teaching
and outraged when Nestorius listened to the complaints of some Alexandrian clergy Cyril had
disciplined.

Thus, late in 428, Cyril opened his attack on Nestorius. He stirred up charges against him and
slandered him in Rome where Patriarch (Pope) Celestine was upset about Nestorius’s welcome of
certain exiles from Rome.

At the General Council of Ephesus (431), called by Emperor Theodosius II, who had until then
supported Nestorius, Cyril got Nestorius deposed before the late arrival of his Syrian supporters.
When the Syrians arrived, under the leadership of John, Patriarch of Antioch, they proceeded to
condemn Cyril and his followers. Finally, the Roman legates arrived and approved Cyril’s action. The
whole affair was disgustingly riddled with power politics. American church historian Williston
Walker called it “one of the most repulsive contests in church history.”

Eventually Theodosius II surrendered to pressure and expelled Nestorius from the capital. He died
around 450, an exile in Egypt. Most of his supporters, however, refused to accept his
excommunication. To this day it remains unclear to what extent Nestorius’s teachings were actually
heretical and to what extent he suffered as a victim of misunderstanding and misrepresentation.

Nestorius’s followers fled to Persia and founded there the Nestorian Church, which soon enjoyed an
active life. A vital monasticism, an eminent theology, and an imposing missionary activity testify
to its strength. Its missionaries penetrated to Malabar, India, and Turkestan. Between 780 and 823
Nestorian believers even made their way to Tibet and central China. At the beginning of the
fourteenth century the Nestorian Church in central Asia could count ten major churches and numerous
native clergymen. Unfortunately, during the bloody persecution of the Muslim conqueror Tamerlane
(1380), this mission was destroyed.

Today the Nestorian Church in the Near East and India still counts about eighty thousand members,
and in America twenty-five thousand.

In his autobiography Nestorius insisted that he did not oppose the use of God-bearer because he
denied the Godhead of Christ but to emphasize that Jesus was born as a genuine human being with body
and soul. His concerns were not unfounded. Nevertheless, even if unfairly, the church has attributed
the name Nestorius to the failure to view the divine and human natures as truly united.

Some recent liberal understandings of Christ are labeled as Nestorian. They reason, if Nestorius
believed that only the will power of the human man Jesus held him in a moral and volitional union
with divine Word, then the difference between Christians and Christ himself is one of degree. Jesus
was more attentive and submissive to God than we are. Jesus was not divine but only a loftier
picture of how close a person can be to God; Jesus is a human model but not a divine savior.

Sadly, some evangelical believers fall sway to the same failure in doctrine; they seek Jesus as a
role model for self-help and ignore the life-giving transformation offered by Emmanuel: God with us.

\section{ESTABLISHING THE BOUNDARIES OF TRUTH}

Soon after the Council at Ephesus (431), the third heresy arose. Eutyches, the spiritual leader of a
monastery near Constantinople, defended the one nature in Christ (monophysitism). He combined the
two natures so intimately that the human nature appeared completely absorbed by the divine one. Just
“as a drop of honey, which falls into the sea, dissolves in it,” so the human nature in Christ is
lost in the divine. Thus Eutyches denied the central prerequisite for the mystery of Christ and his
mission as Savior and Redeemer. The whole Christian doctrine of redemption was in danger.

Patriarch Flavian of Constantinople called Eutyches before a synod and when he refused to recant,
Flavian condemned him as a heretic. Eutyches, however, found support in Dioscorus, Patriarch of
Alexandria, who followed Cyril’s ideas. At Dioscorus’s request, Emperor Theodosius II once again
summoned an imperial council. It assembled under Dioscorus’s leadership in Ephesus (449) and
rehabilitated Eutyches, even though it was not recognized by the rest of the church. Pope Leo I
(440–461) called it the “Robber Council.” He supported the patriarch of Constantinople and asked the
emperor for a new council. The successor of Theodosius, Emperor Marcian (450–457), granted the
request and in 451 called the fourth General Council of Chalcedon.

At this town not far from Constantinople, nearly four hundred bishops gathered and quickly indicted
Dioscorus for his actions at the robber council. Then the assembled fathers, despite some reluctance
to supplement Nicaea, set forth a new definition:

\begin{blockquote}
We all with one voice confess our Lord Jesus Christ one and the same Son, at once complete in
Godhead and complete in manhood, truly God and truly man,~.~.~. acknowledged in two natures, without
confusion, without change, without division, or without separation; the distinction of natures being
in no way abolished because of the union, but rather the characteristic property of each nature
being preserved, and coming together to form one person.
\end{blockquote}

So against Arius the church affirmed that Jesus was truly God, and against Apollinaris that he was
truly man. Against Eutyches it confessed that Jesus’ deity and humanity were not changed into
something else, and against Nestorius that Jesus was not divided but was one person.

From that date forward most Christians in Catholicism, Protestantism, and Orthodoxy looked to
Chalcedon for the foundation of the doctrine of salvation, a unique God-man, Jesus Christ.

Numbers of Christians in the Near East, however, rejected the work of Chalcedon. They held that
instead of the divine and human natures joining to form one person in Jesus, he possessed but one
nature in which divine life and human were indistinguishable. This monophysite (one nature) teaching
was an important factor contributing to the breaking away of the Monophysite Churches from the rest
of Eastern Orthodoxy. Coupled with the decline of the Byzantine power in the outlying areas of the
Eastern Empire, monophysite doctrine led to the Coptic Church, the largest Christian body in Egypt
today, with a related church in Ethiopia, and the so-called Jacobite Church of Syria, which has most
of its adherents in South India.

Obviously Chalcedon did not solve the problem of how deity can unite with humanity in a single
person. At the human level the problem resists explanation. The Bible regards the Event as
absolutely unique. The merit of the Chalcedonian statement lies in the boundaries it established. In
effect, it erected a fence and said, “Within this lies the mystery of the God-man.” Fifteen hundred
years after the event we may wish for more understandable terms, but we dare not say less than the
church said then.

\readinghead{Suggestions for Further Reading}
\begin{readinglist}
\entry Bevan, Edwyn. Christianity. New York: Henry Holt \& Co., 1932.
\entry *Ferguson, Everett. Church History, Volume 1: From Christ to the Pre-Reformation. Grand
  Rapids: Zondervan, 2005.
\entry Hardy, Edward R., ed. Christology of the Later Fathers. Philadelphia: The Westminster Press,
  1954.
\entry Kelly, J. N. D. Early Christian Doctrines, 5th ed. New York: Harper, 1978. Continuum, 2000.
\entry Prestige, G. L. Fathers and Heretics. London: S.P.C.K., 1963.
\entry Wand, J. W. C. Four Great Heresies. London: A. R. Mowbray, 1955.
\end{readinglist}
